% =============================================================
% การพัฒนาเกม
% โครงงาน: การออกแบบและพัฒนาเกมเดินเรือด้วยระบบสั่งการด้วยเสียงเพื่อส่งเสริมทักษะการสื่อสาร
% =============================================================

\section{การพัฒนาเกม}

การพัฒนาเกมดำเนินการบนเอนจิน Unity และใช้ภาษา C\# เป็นหลัก โดยแบ่งการพัฒนาออกเป็นระบบย่อยและการบูรณาการเทคโนโลยีปัญญาประดิษฐ์เข้ากับกลไกของเกม ดังนี้

\subsection{การพัฒนาระบบรับคำสั่งเสียง (Speech Recognition System)}
ระบบรับคำสั่งเสียงพัฒนาโดยเชื่อมต่อไมโครโฟนเข้ากับโมเดลการรู้จำเสียง Whisper~\cite{radford2023} ซึ่งทำงานภายในเครื่องของผู้เล่นผ่านแพ็กเกจ whisper.unity (ไลบรารี whisper.cpp) ด้วยโมเดล ggml-base.en สำหรับภาษาอังกฤษ ระบบรับเสียงพูดภาษาอังกฤษของผู้เล่นขณะกดปุ่มวิทยุ และแปลงเสียงเป็นข้อความ (Speech-to-Text) หลังผู้เล่นปล่อยปุ่มวิทยุ จากนั้นส่งข้อความที่ได้เข้าสู่ตัวแยกคำสั่ง (Command Parser) เพื่อนำไปสั่งการเรือต่อไป ดังภาพที่ \ref{fig:code_speech} (ดูภาคผนวก ข ภาพที่ \ref{fig:code_stt})

\begin{figure}[!htbp]
\centering
\begin{lstlisting}[language={[Sharp]C},basicstyle=\ttfamily\footnotesize,firstnumber=522]
// Only the Whisper call is guarded here, so a bug in routing is never reported
// to the player as a speech error.
string text;
_transcribing = true;
try
{
    var result = await _whisperManager.GetTextAsync(speech, SAMPLING_RATE, channels);
    text = (result != null && result.Result != null) ? result.Result.Trim() : null;
}
catch (Exception ex)
{
    Debug.LogException(ex);
    _transcribing = false;
    Reject(RadioStatus.Failed, "Radio error. Try again.", "(transcription failed)");
    return;
}
_transcribing = false;
\end{lstlisting}
\vspace{0.3em}
\begin{lstlisting}[language={[Sharp]C},basicstyle=\ttfamily\footnotesize,firstnumber=555]
if (CommandRouter.LooksLikeNoise(text))
{
    LanguageLog.Call();
    LanguageLog.Unclear();
    Reject(RadioStatus.NotUnderstood, "Didn't catch that. Try again, short and clear.", "(nothing intelligible: \"" + text + "\")");
    return;
}
\end{lstlisting}
\vspace{0.3em}
\begin{lstlisting}[language={[Sharp]C},basicstyle=\ttfamily\footnotesize,firstnumber=578]
CommandRouter router = CommandRouter.Instance;
if (router == null)
{
    router = FindAnyObjectByType<CommandRouter>();
}

if (router != null)
{
    router.RouteCommand(text, modeUsed == RadioMode.Conversation);
}
\end{lstlisting}
\caption{รหัสต้นฉบับส่วนการรับคำสั่งเสียงและการแปลงเป็นข้อความ (RadioMic.cs บรรทัด 522--538, 555--561, 578--587)}
\label{fig:code_speech}
\end{figure}

\subsection{การพัฒนาระบบลูกเรือปัญญาประดิษฐ์ (AI Crew System)}
ระบบนี้ใช้บริการ API ของ OpenAI (GPT-4o mini) ช่วยตีความและโต้ตอบ ข้อความคำสั่งจะผ่านตัวแยกคำสั่งภายในเครื่องก่อน หากไม่พบคำสั่งที่ชัดเจนจึงส่งให้ GPT-4o mini จำแนกเจตนา ส่วนการสนทนากับลูกเรือ (ภาพที่ \ref{fig:code_ai_crew}) ส่งบุคลิก สถานะเรือ และบทสนทนาล่าสุดให้ GPT-4o mini สร้างคำตอบ หากยังไม่เข้าใจคำสั่งหรือบริการขัดข้อง ลูกเรือจะขอให้ผู้เล่นพูดซ้ำ (ดูภาคผนวก ข ภาพที่ \ref{fig:code_ai})

\begin{figure}[!htbp]
\centering
\begin{lstlisting}[language={[Sharp]C},basicstyle=\ttfamily\footnotesize,firstnumber=164]
private IEnumerator ConverseCoroutine(RadioChannel crew, string speech, bool conversational)
{
    _busy = true;
    int epoch = RadioEpoch.Current;

    var telemetry = ShipTelemetry.Capture();
    string crewName = CrewVoice.DisplayName(crew);

    List<string> log;
    _history.TryGetValue(crew, out log);
    string prompt = CrewPersona.StatelessPrompt(telemetry, lastFinnReport, crew, log, speech);

    var requestData = new OpenAIRequest(
        Model,
        new List<OpenAIMessage>
        {
            new OpenAIMessage("system", CrewPersona.SystemPrompt(crew)),
            new OpenAIMessage("user", prompt)
        },
        Temperature,
        TopP,
        RepeatPenalty,
        MaxTokens
    );

    string body = JsonUtility.ToJson(requestData);
\end{lstlisting}
\caption{รหัสต้นฉบับส่วนการสร้างคำขอสนทนาของลูกเรือปัญญาประดิษฐ์ (OpenAIClient.cs บรรทัด 164--189)}
\label{fig:code_ai_crew}
\end{figure}

\subsection{การพัฒนาระบบฟิสิกส์และการควบคุมเรือ (Ship Physics)}
การควบคุมเรือเขียนด้วยภาษา C\# เพื่อจำลองการเคลื่อนที่ของเรือขนาดใหญ่ โดยคำนวณความเฉื่อย (Inertia) อัตราการหันที่เพิ่มขึ้นทีละน้อย (Yaw Acceleration) และการไถลด้านข้าง (Lateral Drift) ทำให้เรือไม่ตอบสนองต่อการเลี้ยวหรือการหยุดเรือในทันทีที่ได้รับคำสั่ง ดังภาพที่ \ref{fig:code_physics} (ดูภาคผนวก ข ภาพที่ \ref{fig:code_engine})

\begin{figure}[!htbp]
\centering
\begin{lstlisting}[language={[Sharp]C},basicstyle=\ttfamily\footnotesize,firstnumber=151]
private void UpdateRotation(float rudder)
{
    float speed = _engine.Speed;

    // Rudder bite ramps from a slow standstill pivot up to full authority,
    // and reverses when the ship is making way astern.
    float bite = Mathf.Lerp(PivotAuthority, 1f, Mathf.Clamp01(Mathf.Abs(speed) / FullAuthoritySpeed));
    float cargo = GameSession.Instance != null ? GameSession.Instance.CargoTurnFactor : 1f;
    // The rudder upgrade turns the ship itself faster, not just the rudder blade.
    float helm = GameSession.Instance != null ? GameSession.Instance.HelmTurnFactor : 1f;
    float target = rudder * TurningPower * helm * cargo * bite * Sternward;

    _turnRate = Mathf.MoveTowards(_turnRate, target, YawAccel * helm * Time.deltaTime);
    if (Mathf.Abs(_turnRate) < 0.01f) return;

    transform.Rotate(0f, _turnRate * Time.deltaTime, 0f, Space.World);
}

private void UpdateTranslation(float rudder)
{
    if (_knockbackVelocity.sqrMagnitude > 0.001f)
    {
        transform.position += _knockbackVelocity * Time.deltaTime;
        _knockbackVelocity = Vector3.Lerp(_knockbackVelocity, Vector3.zero, Time.deltaTime * 4f);
    }

    float speed = _engine.Speed;
    if (Mathf.Abs(speed) < 0.01f) return;

    // Slide away from the turn so the stern kicks out. The turn reverses when going
    // astern, so the slide does too. Flattened to the water plane.
    Vector3 drift = transform.right * (-rudder * Sternward * LateralDrift * Mathf.Abs(speed));
    drift.y = 0f;

    Vector3 forward = transform.forward;
    forward.y = 0f;
    transform.position += (forward.normalized * speed + drift) * Time.deltaTime;
}
\end{lstlisting}
\caption{รหัสต้นฉบับส่วนระบบฟิสิกส์และการจำลองการเดินเรือ (ShipMovement.cs บรรทัด 151--188)}
\label{fig:code_physics}
\end{figure}

\subsection{ระยะการพัฒนาเกม (Development Phases)}
การพัฒนาเกมตัวต้นแบบแบ่งออกเป็น 4 ระยะหลัก เพื่อให้การทำงานเป็นระบบและทดสอบได้ในแต่ละขั้นตอน ดังนี้

\begin{enumerate}[label=\arabic*)]
    \item \textbf{ระยะที่ 1 (Prototype):} พัฒนาระบบพื้นฐาน โดยเน้นการเชื่อมต่อไมโครโฟนเข้ากับโมเดล Whisper และทดสอบการแปลงเสียงเป็นข้อความแล้วส่งเข้าตัวแยกคำสั่ง เพื่อขยับวัตถุเรือเบื้องต้น
    \item \textbf{ระยะที่ 2 (Alpha):} พัฒนาระบบสภาพแวดล้อมทางทะเล การสุ่มสร้างสิ่งกีดขวาง (Obstacle Generation) และเพิ่มเสียงพูดของลูกเรือปัญญาประดิษฐ์
    \item \textbf{ระยะที่ 3 (Beta):} พัฒนาระบบ UI/UX ให้สมบูรณ์ ปรับจูนความยากง่าย (เช่น ความเร็วของเรือ ระยะห่างของสิ่งกีดขวาง) และความแม่นยำ (Tolerance) ในการรับคำสั่งเสียง
    \item \textbf{ระยะที่ 4 (Release Candidate):} ตรวจสอบและแก้ไขข้อผิดพลาด (Bug Fixing) และเตรียมแพ็กเกจเกมสำหรับนำไปทดสอบกับผู้ใช้งานจริง
\end{enumerate}